% ===== Fuentes (B25-26, actualizadas) =====
\begin{frame}{Fuentes (1/2)}
  \vspace{0.2em}
  {\fontsize{6.6}{8.4}\selectfont\raggedright
  \textbf{Demografía.} Naciones Unidas, DESA, División de Población, \emph{World Population Prospects 2024}, variante media, México; grupos de edad agregados a partir del archivo oficial por grupos quinquenales (cálculos propios). La revisión 2026 fue pospuesta a 2027. INEGI, Encuesta Nacional de la Dinámica Demográfica 2023. INEGI, Estadística de Nacimientos Registrados 2024. Fernández-Villaverde, J. (2023), \emph{The Demographic Future of Humanity: Facts and Consequences}.\par\vspace{0.6em}
  \textbf{Ciclo de vida y efecto de la demografía sobre el gasto.} National Transfer Accounts (ntaccounts.org), perfiles por edad de México, 2004 y 2014. Gasto con perfiles fijos: cálculo propio con el perfil de 2004 y una proyección de población propia con datos del INEGI; resultados en \% del ingreso laboral total, no del PIB.\par\vspace{0.6em}
  \textbf{Presupuesto 2027, salud y pensiones.} PEF 2027: dictamen de la Cámara de Diputados; gasto neto total de 10.64 billones de pesos. CIEP (2026), \emph{Implicaciones del Paquete Económico 2027}, 11 de septiembre de 2026: gasto en salud, gasto per cápita por subsistema, gasto de bolsillo, inflación de salud, carga de enfermedad (IHME, 2021; OMS, 2024), prima de equilibrio del IMSS (IMSS, 2024), pensiones, Miscelánea e ingresos petroleros; sus cifras corresponden al proyecto de presupuesto. Meta de 6\% del PIB: OPS (2014), \emph{Estrategia para el acceso universal a la salud y la cobertura universal de salud}.\par\vspace{0.6em}
  \textbf{Espacio disponible.} Villarreal, \emph{Contabilidad de la Muerte}, actualización con el Informe de Finanzas Públicas y Deuda Pública enero--septiembre 2026 (SHCP) y los Criterios Generales de Política Económica 2027. PIB nominal de 36,500 mmdp; el supuesto favorece a la cuenta oficial. CIEP (2026) estima un espacio fiscal de 1.6\% del PIB con una definición distinta.\par\vspace{0.6em}
  \textbf{Deuda y reforma al ISR.} LIF 2027, aprobada por el Senado el 27 de octubre de 2026, arts. 1o, 2o y 4o. SHCP, Criterios Generales de Política Económica 2027. Congreso de la Unión, minuta de reforma a la LISR 2027.\par}
\end{frame}

\begin{frame}{Fuentes (2/2)}
  \vspace{0.2em}
  {\fontsize{6.6}{8.4}\selectfont\raggedright
  \textbf{Reforma fiscal clásica.} Estimación propia del rendimiento de un paquete de reforma clásica (IVA, ISR y gastos fiscales), alrededor de 1.5 puntos del PIB.\par\vspace{0.6em}
  \textbf{Trampa de la formalización.} Fietz, Joubert, Ñopo, Ocampo, Packard, Posadas y Rodriguez Chamussy (2025), \emph{(In)Formalizing Jobs in Latin America and the Caribbean: Taxes, Benefits, and Labor Market Incentives}, Banco Mundial, cap. 5. Uruguay: Banco de Previsión Social, resoluciones de directorio 2021--2026; Ministerio de Economía y Finanzas, Rendición de Cuentas 2023.\par\vspace{0.6em}
  \textbf{Cuota fija.} Ley del Seguro Social, art. 106, fr. I. IMSS, \emph{Informe al Ejecutivo Federal y al Congreso de la Unión sobre la situación financiera y los riesgos del IMSS 2025--2026}: cuota fija patronal de 20.40\% de la UMA. INEGI, UMA 2026: 117.31 pesos diarios. CONASAMI, salario mínimo general 2026: 315.04 pesos diarios.\par\vspace{0.6em}
  \textbf{Infonavit.} Ejercicio propio, preliminar, sobre tasas legales. No es una propuesta del IMSS ni del Gobierno Federal. Antecedente: Villarreal, H. (2022), «El Infonavit, ¿qué sigue?», \emph{Arena Pública}, 20 de diciembre de 2022.\par\vspace{0.6em}
  \textbf{Modalidad 40.} Acuerdo ACDO.AS2.HCT.101026/302.P.DG del Consejo Técnico del IMSS, \emph{Diario Oficial de la Federación}, 21 de octubre de 2026. Ley del Seguro Social de 1973.\par\vspace{0.6em}
  \textbf{Modelo de aseguramiento y metas de financiamiento.} Villarreal, H. (2024), «¿Se puede reformar la seguridad social? Parte II», \emph{Arena Pública}, 10 de septiembre de 2024. Dilema de Gruber: Jonathan Gruber (MIT).\par\vspace{0.6em}
  \textbf{Arquitectura institucional.} Diario Oficial de la Federación: 29-XI-2019 (Insabi), 31-VIII-2022 (IMSS-Bienestar), 29-V-2023 (reforma a la Ley General de Salud).\par\vspace{0.6em}
  \textbf{Comunidad.} Documento de presentación de la iniciativa; reunión de trabajo en San José, Costa Rica, 5--7 de octubre de 2026.\par}
\end{frame}
